\chapter{线性回归与正则化}
\chaptergoal{理解线性模型的预测表达和平方误差目标；使用 sklearn 拟合单特征与多特征回归；区分求解器和模型本身；通过残差、特征变换与正则化理解模型偏差和稳定性；完成一个具有专业单位解释的回归实践。}

线性回归是统计学与机器学习之间的重要桥梁。物理学生可能已经见过直线拟合、最小二乘和误差传播，但在机器学习中，除了求出参数，还要关注对新数据的预测、变量处理、交叉验证和模型适用范围。本章让同一套回归数据贯穿多个概念，并在最后将流程迁移到天文或实验室测量。
\section{从单特征回归到多特征回归}

\label{sec:housing-intro}
\textbf{线性回归（linear regression）}用输入特征的加权和加上一个常数来预测数值目标。这里的“权重”就是各输入前面的系数；常数项也叫截距。训练的任务是从已知输入与目标的记录中确定这些系数。第3章已经用一条直线说明训练机制，本章把它扩展为可以实际使用的回归模型。

先看一个人工规定的两特征模型：$\hat y=1+2x_1-0.5x_2$。若某条记录的输入为 $(x_1,x_2)=(3,4)$，预测为 $1+2\times3-0.5\times4=5$。此处两个特征和目标均无物理单位，只用于说明计算。模型使用的是同一套系数；换一条记录时改变输入，而不是为每条记录重新拟合一组系数。

只有一个输入时，线性回归为 $\hat y=wx+b$；有 $p$ 个输入时，写成
\begin{equation}
\hat y_i=b+\sum_{j=1}^{p}w_jx_{ij}=\mathbf x_i^{\mathsf T}\mathbf w+b.
\end{equation}
$w_j$ 表示在其他输入保持不变时，该输入变化一个单位对应的线性预测变化量。这个数学解释并不自动等同于因果效应：高度相关的输入、遗漏变量和选择偏差都可能改变参数含义。

\subsection*{认识加州住房数据：一条记录是一个区域}
本章使用\textbf{加州住房数据集（California Housing）}。它根据1990年美国加州的人口普查资料整理，包含20\,640条区域记录和8个输入特征。每条记录对应一个人口普查街区组，即由若干住户组成的统计区域，而不是一套房屋。要预测的目标是该区域住房价值的中位数，单位为十万美元。因此目标值为2.5表示区域住房价值中位数为25万美元，并不表示本次程序输出了某套房屋的实际成交价格。\textsuperscript{[10]}

\begin{table}[htbp]\centering\small
\begin{tabularx}{\linewidth}{lX}
\toprule 字段名 & 记录的含义 \\\midrule
\texttt{MedInc} & 区域收入中位数，沿用数据源的缩放口径，不与目标的十万美元单位混同 \\
\texttt{HouseAge} & 区域房龄中位数，单位为年 \\
\texttt{AveRooms} & 每户平均房间数 \\
\texttt{AveBedrms} & 每户平均卧室数 \\
\texttt{Population} & 区域人口数 \\
\texttt{AveOccup} & 每户平均居住人数 \\
\texttt{Latitude}、\texttt{Longitude} & 区域纬度和经度，单位为度 \\
\texttt{MedHouseVal} & 目标：区域住房价值中位数，单位为十万美元 \\
\bottomrule
\end{tabularx}
\caption{加州住房数据的输入与目标。最后一行是预测目标，不属于8个输入特征。字段含义依据数据源说明。}
\end{table}

例如，第一条区域记录的房龄中位数为41年、人口为322，目标值为4.526（即45.26万美元）。这几个值只是帮助我们读懂一行记录；正式模型使用哪些输入，应在实验开始时明确。上述历史样本不代表当前市场价格，也不能直接检验模型对其他年代的预测能力。\textsuperscript{[10]}

下面把8个输入放在\texttt{X}中，把目标放在\texttt{y}中。先划出保留数据，暂不使用；在开发部分拟合一个线性模型。代码中的\texttt{fit}根据配对记录求参数，\texttt{predict}则用求得的参数计算预测。

\begin{lstlisting}[caption={加州住房数据的读取与基线回归（首次运行需联网）}]
from sklearn.datasets import fetch_california_housing
from sklearn.model_selection import train_test_split
from sklearn.linear_model import LinearRegression
from sklearn.metrics import mean_absolute_error, root_mean_squared_error

housing = fetch_california_housing(as_frame=True)
X, y = housing.data, housing.target
X_dev, X_test, y_dev, y_test = train_test_split(
    X, y, test_size=0.2, random_state=42
)
reg = LinearRegression().fit(X_dev, y_dev)
print("train MAE", mean_absolute_error(y_dev, reg.predict(X_dev)))
# Reserve X_test/y_test until the workflow is fixed.
\end{lstlisting}

首次离线运行若未缓存数据，会报数据不可获取的错误。教材工程附有单独的数据准备脚本，教师应在联网机器预先获取并记录软件版本。与其用随机生成的数据偷偷代替住房数据，不如明确说明下载依赖，并把模拟数据单独标注为机制演示。

\subsection{线性模型实际上在假设什么}
线性回归不是要求所有观测点恰好落在同一条直线上，而是用一个线性函数描述在所选输入下的预测趋势。可以把观测写成$y_i=b+\mathbf x_i^{\mathsf T}\mathbf w+\varepsilon_i$，其中$\varepsilon_i$概括没有由这些输入解释的差异。差异可能含有测量误差、遗漏变量和模型近似误差，不能一律解释成仪器噪声。

在平方误差评价下，若允许任意函数并且相关期望存在，最佳总体预测是条件均值$\mathbb E[Y\mid X=x]$。线性回归把这个搜索限制在所选线性函数族内。这个判断可以从$\mathbb E[(Y-c)^2\mid X=x]=\operatorname{Var}(Y\mid X=x)+(\mathbb E[Y\mid X=x]-c)^2$看出：第一项不随$c$改变，第二项在$c$等于条件均值时最小。真实条件均值不一定线性，所以即使最小二乘求解很准确，模型形式仍可能不足。

\subsection{系数的单位与交互作用}
若$y$以伏特计、$x_1$以摄氏度计，$w_1$的单位就是伏特每摄氏度；截距与$y$同单位。每列尺度不同并不妨碍最小二乘拟合，但会影响系数的数值大小。比较“哪一系数绝对值更大”之前，先检查单位，不能直接把大系数等同于更重要的物理量。

线性模型还可以包含交互特征$x_1x_2$。例如传感器电压响应随温度变化，可写成$\hat y=b+w_1x+w_2T+w_3xT$。此时对输入电压的局部斜率是$w_1+w_3T$，不再是固定的$w_1$。模型对参数仍然线性，但变量的解释必须结合其他项。第17章的模拟标定项目将比较是否加入这个温度依赖关系。

一个变量被写进模型，不意味着其系数一定能由现有数据准确确定。若所有实验都在同一温度进行，温度效应与截距难以分离。采样范围和实验设计决定可辨识信息；算法名称不能补救缺乏变化的输入。

\section{平方损失与最小二乘}
从第3章的残差 $e_i=\hat y_i-y_i$ 出发，普通最小二乘最小化
\begin{equation}
L(\mathbf w,b)=\sum_{i=1}^n(y_i-\mathbf x_i^{\mathsf T}\mathbf w-b)^2.
\end{equation}
与半均方误差相比，这里只相差一个不依赖参数的正比例因子，不改变无正则化情况下的最优点，但会影响某些带正则化目标中系数的数值定义。平方损失倾向于重视较大残差，因此数据中少数异常测量可能对拟合产生较大影响。

\subsection*{算一次三点直线拟合}
考虑三条人工记录 $(x,y)=(0,1),(1,2),(2,2)$。它们没有专业物理含义，用于核对最小二乘计算。输入均值为 $\bar x=1$，目标均值为 $\bar y=5/3$。对于有截距的一元直线，最小二乘系数可写成
\[
w=\frac{\sum_i(x_i-\bar x)(y_i-\bar y)}{\sum_i(x_i-\bar x)^2}
=\frac{(-1)(-2/3)+0(1/3)+1(1/3)}{1+0+1}=\frac12,
\qquad b=\bar y-w\bar x=\frac76.
\]
这里分子描述输入与目标共同变化的程度，分母描述输入自身的变化程度；若所有输入相同，分母为零，便不能由这条公式确定斜率。代入三个输入，预测分别是 $7/6,5/3,13/6$，残差 $\hat y-y$ 分别是 $1/6,-1/3,1/6$，平方和为 $1/6$。作为比较，始终预测目标均值 $5/3$ 的常数模型，其残差平方和为 $2/3$。直线利用输入信息减少了这组三点的训练误差；泛化仍需另行评价。

\begin{lstlisting}[caption={用程序核对三点最小二乘结果}]
import numpy as np
from sklearn.linear_model import LinearRegression
x = np.array([0., 1., 2.]).reshape(-1, 1)
y = np.array([1., 2., 2.])
fitted = LinearRegression().fit(x, y)
print(fitted.coef_[0], fitted.intercept_)
print(np.sum((fitted.predict(x) - y) ** 2))
\end{lstlisting}
这里\texttt{reshape(-1, 1)}把三个数排成三行一列的特征表，\texttt{coef\_}和\texttt{intercept\_}保存已经学习的斜率与截距。应得到约 $0.5$、$1.16667$ 和 $0.16667$。\texttt{LinearRegression}的常规稠密实现使用最小二乘求解器，并不是运行第3章写出的梯度下降循环。\textsuperscript{[6]}

\subsection{从最小化目标推到直线公式}
令$S(w,b)=\sum_i(y_i-wx_i-b)^2$。对$b$求导并令其为零，可得$\sum_i(y_i-wx_i-b)=0$，即$b=\bar y-w\bar x$。将其代回目标，相当于对中心化数据拟合过原点的直线。再对$w$求导，得到
\[
\sum_i(x_i-\bar x)(y_i-\bar y-w(x_i-\bar x))=0,
\]
只要$\sum_i(x_i-\bar x)^2>0$，就得到前面的斜率公式。这里两步求导不仅给出计算方法，也给出可检查的性质：\textbf{有截距的普通最小二乘，训练残差总和为零；残差与每个已使用输入的训练内积也为零。}

这些性质只在对应目标、训练样本和精确最小二乘条件下成立。带权拟合对应加权正交，带正则化的系数不再满足完全相同的输入正交条件，验证残差均值也不保证为零。可以利用性质检查实现，但不能把它误认为新数据误差必然没有偏差。

\worked{一个完整的矩阵核对}
另取人工三点$x=(1,2,3)$、$y=(2,3,5)$，设每行设计向量为$(1,x_i)$，得到
\[
D=\begin{pmatrix}1&1\\1&2\\1&3\end{pmatrix},\quad
D^{\mathsf T}D=\begin{pmatrix}3&6\\6&14\end{pmatrix},\quad
D^{\mathsf T}y=\begin{pmatrix}10\\23\end{pmatrix}.
\]
令$\boldsymbol\beta=(b,w)^{\mathsf T}$，正规方程$D^{\mathsf T}D\boldsymbol\beta=D^{\mathsf T}y$给出$b=1/3,w=3/2$。预测为$(11/6,10/3,29/6)$，目标减预测为$(1/6,-1/3,1/6)$，残差和为零，乘以$x$后求和也为零。它与前面的三点例子相似，但数据和系数不同，配套脚本分别明确记录，不能混用输出。

\subsection{最小二乘拟合不要求先假定正态噪声}
为了计算平方误差最小的参数，不必先证明残差正态，也不必先规定独立同分布的统计模型；这是一个确定的优化问题。但若进一步把系数解释为无偏估计，或者给出经典标准误、显著性检验与置信区间，就需要额外假设并检查其适用性。

入门教材经常把“能拟合”“能预测”“能作参数推断”混在一起。一个模型可以在验证中预测有用，却不适合按最简单的公式报告系数误差；也可以满足某种线性假设，却因数据量很少而预测波动很大。本书主线强调预测与残差分析，推断性结论需要更完整的统计论证。

\subsection{测量精度不相同时的加权拟合}
若每条参考目标有可靠的标准差$\sigma_i$，且希望按这些精度加权，可以最小化$\sum_i(y_i-\hat y_i)^2/\sigma_i^2$。标准差较小的测量获得较大权重。把每行设计向量和目标同时除以$\sigma_i$，即可把这个目标转成一个普通平方和问题。

这些权重不是为得到好看的结果而任意设置。它们应来自测量模型或明确的任务代价；输入自身也有显著测量误差时，仅对目标加权并不能解决完整的误差变量问题。天文周期与星等案例中，需要分别考虑周期误差、测光误差和对象间本征散布。

\section{线性模型的求解与梯度下降回顾}
同一种线性模型可以有不同求解方法。最小二乘求解器从线性代数关系获得系数；梯度下降则从初始参数出发，沿着降低目标的方向反复更新。\texttt{SGDRegressor} 支持逐步梯度式训练，但它与 \texttt{LinearRegression} 的默认优化目标与配置可能不同。进行公平对照时，需要确认损失、是否惩罚参数、截距处理、特征尺度和停止准则。

\begin{lstlisting}[caption={最小二乘与无惩罚随机梯度方法的概念对照}]
from sklearn.linear_model import LinearRegression, SGDRegressor
from sklearn.preprocessing import StandardScaler
from sklearn.pipeline import make_pipeline

least_squares = LinearRegression()
gradient_model = make_pipeline(
    StandardScaler(),
    SGDRegressor(loss="squared_error", penalty=None,
                 random_state=42, max_iter=3000, tol=1e-5)
)
# Coefficient scales differ after standardization.
\end{lstlisting}

由于标准化改变输入表示，两者输出预测可以比较，但直接比较未换算的系数不合理。具体结果还受收敛条件影响，不能因为有限步训练存在差别就断言某一方法的预测模型本质更优。

\subsection{正规方程是数学关系，不意味着程序应该显式求逆}
当设计矩阵列满秩时，最小二乘解可写成$(D^{\mathsf T}D)^{-1}D^{\mathsf T}y$。这个表达便于理解，但数值代码通常不应先构造逆矩阵。QR分解或奇异值分解可以更稳健地求解，尤其当列相关时，先形成$D^{\mathsf T}D$会进一步恶化数值条件。

在NumPy中可用\texttt{np.linalg.lstsq(D,y,rcond=None)}进行核对。它返回参数以及秩等信息。若列不满秩，可能存在多组最小损失参数；特定求解器给出的一个解，并不意味着所有系数都被数据唯一确定。第5.5节会用相关特征说明这一点。

\subsection{比较求解路径时要比较同一个问题}
让直接求解与梯度方法拟合同样的训练集时，应先统一目标和预处理。\texttt{SGDRegressor}默认可包含正则化和学习率调度，若直接与无惩罚最小二乘相比，系数差异可能来自目标不同。即使取消惩罚，有限迭代、停止容差和随机顺序也会造成数值差异。

比较应包括训练目标值、验证预测差异、迭代与时间，而不是只看两组系数是否逐位相等。对于病态数据，参数可以相差较大而预测相近；对于缩放后的数据，先把系数换回原单位才有可比性。每个差异都应找对应原因，不能一概解释为“机器学习方法不稳定”。

\worked{由标准化坐标还原两个系数}
若$z_1=(x_1-10)/2$、$z_2=(x_2-100)/5$，训练模型为$\hat y=3+4z_1-2z_2$，代回得到$\hat y=23+2x_1-0.4x_2$。取$x=(12,105)$，两种写法都给出5。标准化后的系数4与$-2$分别对应每个训练标准差的变化；原系数2与$-0.4$对应原始单位。二者没有矛盾，只是坐标不同。

\section{特征变换与多项式回归}
“线性模型”表示对模型参数线性，而不要求预测一定是原始输入的直线。例如
\begin{equation}
\hat y=b+w_1x+w_2x^2
\end{equation}
仍然对参数 $w_1,w_2$ 线性，却能表示弯曲的输入关系。类似地，对严格为正的输入使用对数，也是一种特征变换；本章末尾介绍的恒星亮度案例将使用对数周期。变换应记录定义、单位及适用范围；不能先看到测试结果再挑选让测试分数最好的变换。

\begin{lstlisting}[caption={多项式特征属于训练流程的一部分}]
from sklearn.preprocessing import PolynomialFeatures, StandardScaler
from sklearn.linear_model import Ridge
from sklearn.pipeline import make_pipeline

poly_model = make_pipeline(
    PolynomialFeatures(degree=2, include_bias=False),
    StandardScaler(),
    Ridge(alpha=10.0)
)
# Fit and compare within development cross-validation.
\end{lstlisting}

多项式次数增加后，可表示的关系更加复杂，但外推时高次项也可能迅速增大。实验应比较训练与验证误差，而不是只比较拟合曲线是否更加贴近训练样本。

\subsection{从物理关系选择基函数}
可以把输入变换写为$\phi_1(x),\ldots,\phi_q(x)$，模型为$\hat y=b+\sum_jw_j\phi_j(x)$。多项式、对数、正弦和物理上给定的无量纲组合，都可以作为候选基函数。\textbf{线性是对待估参数而言，不是对原始坐标图形而言。}

例如周期信号可以先用已知频率的$\sin(\omega t)$和$\cos(\omega t)$作输入，模型学习两个系数；若频率$\omega$也要同时学习，问题就不再是对全部未知量的简单线性最小二乘。这一区别解释了为什么“模型公式看起来像几项相加”还不足以判断求解方法。

对数变换只适用于定义域允许的值。将带单位的周期写成$\log_{10}(P/\mathrm{day})$，比直接写一个有单位量的对数更清楚。更换单位会改变截距，因此记录单位不仅是图注要求，也直接影响参数解释。

\subsection{变换目标与变换输入承担不同任务}
用$\log x$预测原目标$y$，改变的是模型对输入关系的表达；用$x$预测$\log y$，改变的是目标和损失的计量。后一种方法更关注相对偏差，但在对数空间最优的预测，简单反变换后未必是原空间的条件均值。例如指数是非线性函数，通常$\exp(\mathbb E[\log Y\mid X])$不等于$\mathbb E[Y\mid X]$。

因此，报告变换目标的模型时，应同时说明训练损失在哪个空间，最终MAE或RMSE在哪个单位下计算。若用原单位评价，就必须先反变换预测，再与原目标比较。不能把对数RMSE0.1直接描述成0.1伏特。

\subsection{实际对照：次数增加并不等于验证持续改善}
随书离线机制实验使用明确生成的响应
\[
y=1.2+1.5x+0.8\sin(3x)+0.25u+\varepsilon,
\]
其中$x$、$u$和噪声由固定随机数生成，全部标为人工单位。实验为突出多项式表达，暂时只使用$x$；$u$的作用成为未解释差异的一部分。比较1、2、3、5、9和15次多项式，并在每个模型后使用很弱的Ridge惩罚帮助数值稳定。

实际验证RMSE依次约为0.6904、0.6925、0.7046、0.4957、0.2862、0.2928。前三种变化没有形成持续改善，9次在本次验证中更好，15次训练误差略降而验证略差。这个结果回答的是有限候选在这一固定生成样本上的表现，不是“9次多项式总是最佳”。没有把保留测试数据用于选次数。

\fig{v3_ch05_degree.pdf}{人工非线性响应上的多项式比较。训练与验证均使用同一目标单位，模型仅使用$x$，未使用$u$。曲线反映表达能力、有限样本及惩罚共同作用。}

\section{相关特征与系数解释}

\textbf{相关特征}是指在当前数据中一起变化、携带部分重叠信息的输入。设一张表同时记录长度的厘米值 $x_1$ 和毫米值 $x_2=10x_1$。模型中的 $w_1x_1+w_2x_2$ 等于 $(w_1+10w_2)x_1$。系数 $(w_1,w_2)=(2,0)$ 与 $(0,0.2)$ 给出完全相同的预测。数据可以约束它们的组合，却不能仅凭预测误差区分这两组系数。现实中并非总是精确倍数，但接近的冗余也会使系数解释变得不稳定。

若两个输入高度相关，可能存在多组系数产生近似相同的预测。回归结果在新样本上可能稳定，但单个系数对训练样本或预处理改变很敏感。解释系数时应说明其他变量、单位和样本范围。标准化后的系数与原单位下的系数不应直接混用。

例如，住房数据中某些人口统计特征与地理变量存在相关结构。若模型包含经纬度，系数并不表示“把房屋移动一度”会产生何种因果变化；它首先刻画的是当前样本与当前模型中的线性关联。对于科学问题，特征重要性和物理因果机制要分开论证。

\subsection{小数据变化为什么会造成大系数变化}
若$x_2$几乎等于$x_1$，预测主要由$w_1+w_2$控制。增加$w_1$、减少$w_2$可能几乎不改变预测，但会显著改变单个系数。设计矩阵在这种方向上包含很少信息，目标曲面相应较平缓。

随书实验构造两个高度相关输入，多次重采样训练普通最小二乘和Ridge。绘图显示两种系数怎样共同变化，并保存系数之和。应同时检查单个系数和组合预测：只看系数抖动可能夸大预测的不稳定，只看预测稳定又可能掩盖单个物理系数无法可靠解释的问题。

\fig{v3_ch05_collinearity.pdf}{两个高度相关人工特征的重采样系数。不同点表示不同重采样，不是不同物理测量值；Ridge通过惩罚选择较小的系数组合。}

\subsection{模型中的“保持其他变量不变”可能缺少数据支持}
数学上，$w_j$描述其他输入固定时$x_j$变化一单位的预测差。但若数据中$x_1$和$x_2$几乎总是一起变化，固定$x_2$而大幅改变$x_1$可能落到从未观测的组合。模型仍然会给出一个数，却不代表数据已经支持这种干预解释。

在专业分析中，先画出相关变量的联合覆盖，再决定系数解释范围。为了解释某个效应而收集具有独立变化的实验条件，往往比更换回归算法更重要。这是数据建模与实验设计之间的联系。

\section{岭回归与 Lasso}

\textbf{岭回归（Ridge）}和\textbf{Lasso 回归}都在线性回归中加入正则化。正则化是在拟合误差之外增加对模型复杂程度的约束或代价；这里具体体现为惩罚较大的系数。这样，即使两组系数的预测误差相近，训练目标也可能偏向系数较小的一组。惩罚平方和通常称为 $L_2$ 惩罚，惩罚绝对值之和称为 $L_1$ 惩罚；下式中的求和正是这两个名称的含义。

线性回归在多个输入相关、样本有限时可能不稳定。正则化给大系数增加代价，常见目标是
\begin{align}
L_{\mathrm{Ridge}}&=\sum_i e_i^2+\alpha\sum_jw_j^2,\\
L_{\mathrm{Lasso}}&=\frac{1}{2n}\sum_i e_i^2+\alpha\sum_j|w_j|.
\end{align}
这两行使用了与 sklearn 各自实现文档相适配的常见写法；不同模型和教材对损失常数、$n$ 与 $\alpha$ 的组合可能采取不同约定，因此比较不同包的 $\alpha$ 数字不能忽略其目标函数定义。截距通常不受同样的惩罚。

Ridge 倾向于连续收缩系数，Lasso 能使部分系数变为零。非零系数集合不自动等于真实物理变量集合，尤其在相关特征中，Lasso 可能在几乎等价的输入间选择其中一个。实验应该使用训练折内标准化，并在交叉验证中确定 $\alpha$。

一个单参数人工例子能直接显示二者的区别。设无惩罚损失是 $\tfrac12(w-2)^2$，最优系数为2。加入 $\tfrac\lambda2w^2$ 后，令导数为零得到 $w=2/(1+\lambda)$；当 $\lambda=1$ 时，系数缩为1。若改加 $\lambda|w|$，最优值为 $\max(2-\lambda,0)$：$\lambda=1$时为1，$\lambda=2$时恰好为0。这个例子解释“收缩”与“置零”，其中 $\lambda$ 的定义只属于此例，不应直接当作程序中所有模型的\texttt{alpha}。

\subsection*{离线辅助数据：糖尿病进展记录}
为使未下载住房数据的读者也能观察系数随惩罚强度变化的图，本节另用\textbf{糖尿病进展数据集（Diabetes）}制作一张辅助图。它由442名患者的记录组成，每人有10个基线输入，目标是基线一年后的疾病进展定量指标。输入涉及年龄、性别、身体质量指数、平均血压和6项血清测量；\texttt{load\_diabetes()}默认返回按其文档约定中心化和缩放后的输入，目标不是疾病类别。\textsuperscript{[29]}

该数据只承担离线回归演示，不是住房数据的另一名称，也不是本书的医学研究项目。它的目标单位与住房价值不同，二者的误差数字不能直接比较。随书实验明确标出使用哪一数据，不会在下载失败时悄悄替换数据。


\fig{ch05_ridge_path.pdf}{用 sklearn 随附的 Diabetes 数值数据演示 Ridge 系数如何随正则化强度变化。这是优化行为的辅助示意，不是本章的住房主案例。}

\subsection{正则化改变的不只是数值求解}
在线性最小二乘中，Ridge相当于在系数方向增加曲率，可以缓解一些病态问题；同时，它也改变了最优点，允许训练拟合稍差以换取较小系数。不能把正则化仅描述成一个使程序更快的技巧。其验证收益来自对有限样本的限制是否与真实关系相容。

为了看清尺度依赖，设$x'=1000x$。同一预测可用$w'=w/1000$表示；若仍施加数值相同的$\alpha(w')^2$，这一方向受到的惩罚变弱一百万倍。因此，在有惩罚的模型中，单位变换可能改变最终预测，而无惩罚的精确最小二乘预测通常可通过系数换算保持不变。训练折内标准化是常见办法，但实际物理约束也可以采用事先规定的尺度。

\subsection{Lasso为什么能够得到恰好为零的系数}
对单参数例子$J(w)=\tfrac12(w-a)^2+\lambda|w|$，在$w>0$时导数为$w-a+\lambda$，在$w<0$时为$w-a-\lambda$。当$|a|>\lambda$时，解沿原来的符号缩小$\lambda$；当$|a|\le\lambda$时，两侧的斜率使最小点位于零。合起来就是
\[
w^*=\operatorname{sign}(a)\max(|a|-\lambda,0).
\]
这个分段结果称为软阈值形式。它解释了尖角为什么允许精确置零，而Ridge的光滑惩罚通常只连续缩小系数。多特征相关情形更复杂，不能对原始最小二乘系数逐个简单套这条公式；但单参数计算足以建立稀疏性的直觉。

\subsection{验证选择与变量选择的区别}
选择使验证误差较小的惩罚强度，是预测模型选择；把最终非零变量解释为真实机制，是更强的科学主张。相关输入可能互相替代，样本稍变就换一个非零变量。可以报告不同重采样中变量入选的频率，检查预测是否随选择改变，但这些仍不自动证明因果关系。

实际实验应同时报告惩罚强度、标准化方法、非零系数数目和验证误差。一个只剩两列的模型与保留全部输入的模型分数接近时，简洁性可能有价值；但应确认这两列在部署时能够稳定测量，而不是只因数据源中恰好保存得好。

\section{预测误差、残差与适用范围}
回归模型应提交预测与真实目标的散点图、残差与预测值关系、残差分组统计及 MAE、RMSE 等指标。残差随输入呈系统曲线可能说明模型形式过于简单；在高目标值处偏差增大可能反映样本不足、异方差或目标边界条件。分析时不能仅凭图形确定原因，还要结合数据来源、测量方法和模型假设。

正式评价时关注独立测试的整体指标，同时在开发范围内先确定哪些子群体值得诊断。若某些记录属于不同地区、仪器或物理工况，需要报告不同条件下的表现，而不只给出单一平均分。将模型用于训练范围之外时，应明确这是外推，而不是已验证插值。这里“插值”指在已有输入覆盖范围内估计，“外推”指延伸到覆盖范围之外；二者对证据的要求不同。

\subsection{四类残差模式与下一步动作}
残差整体偏正或偏负，提示系统偏移；残差随输入呈弯曲，提示线性表达可能不足；残差散布随预测增大，提示误差尺度可能变化；某些单独对象偏离很大，提示需要检查测量、标签、覆盖范围或稀有物理状态。每一种图形都只是线索，下一步应通过数据字段和对照模型检验解释。

例如弯曲残差可以用事先有意义的二次项或对数输入进行比较，检查开发验证是否改善；不能只在看到测试残差后不断增加项，再把该测试分数称为最终独立表现。异常点也不应仅因降低拟合分数而删除，应记录排除依据和排除前后变化。

\subsection{外推能力来自模型假设，不来自边缘分数}
直线模型在训练范围外仍沿直线延伸，多项式模型可能迅速增长，树模型常落入某个边界叶而输出常数。三者都会给出数字，但数字来源不同。输入位于各列范围之内，也不保证处于训练样本实际覆盖的联合区域。例如温度和压力各自都在范围内，它们的某种组合却可能从未出现。

专业报告应明确要求是同条件下的新对象、不同工况还是超出范围的预测。可以把某些极端工况完整留出，形成更接近外推的评价，但这会改变问题，不宜与随机划分分数混为一谈。

\section{专业回归实践}
\subsection*{天文方向：造父变星周期--星等关系}

\label{sec:ogle-intro}
\textbf{造父变星}是一类亮度随时间近似周期变化的脉动恒星。这里要利用的量是完成一个亮度循环所需的时间，即周期 $P$，以及在选定观测波段中的平均亮度。天文学常用\textbf{星等}表示亮度：在同一标定体系下，数值越小表示越亮；“视星等”描述从观测者处看到的亮度，而不是直接描述恒星本身的总发光能力。

本案例的数据来自\textbf{光学引力透镜实验（Optical Gravitational Lensing Experiment，OGLE）}的变星目录。目录给出大麦哲伦云方向经典造父变星的编号、脉动模式、周期和平均测光量。其中 I 波段表示选定的光学观测波段；所用目录字段是强度平均 I 波段星等，单位为mag。每条记录对应一颗恒星，而不对应某一次曝光。使用同一地区、同一脉动模式，可以先减少样本关系中的部分差异，但并不消除全部物理影响。\textsuperscript{[13]}

选定 OGLE 同一目标区域、同一脉动模式的经典造父变星，读取周期 $P$（天）与平均 I 波段视星等 $m_I$，在数据质量检查后拟合 $m_I=a\log_{10}(P/\mathrm{day})+b$。先用简单的线性基线，检查残差与周期及星等的关系。这个练习研究的是选定样本中的\textbf{周期--视星等关系}，不能把它直接称为不依赖消光和距离的绝对光度标定。需记录筛选、来源版本和误差处理。

\subsection*{物理方向：风洞声学实验}

\label{sec:airfoil-intro}
\textbf{翼型自噪声数据集（Airfoil Self-Noise）}来自风洞中的翼型声学测量。翼型是机翼横截面的形状；风洞通过受控气流研究它周围的流动及发声。数据共有1503条记录，学习任务是根据实验条件预测声压级，而不是识别某一架飞机。\textsuperscript{[16]}

五个输入依次为频率（Hz）、攻角（度，翼型与来流的相对角度）、弦长（m，翼型前后缘之间的特征长度）、自由来流速度（m/s）和吸力侧位移厚度（m，用于描述近壁流动影响的量）。目标是数据源整理的声压级（dB）。一行描述一组给定条件下的测量；同一翼型或工况的重复测量不应自动当作独立实验。原始文件和字段顺序需依数据字典读取，不能把最后一列误作输入。

UCI Airfoil Self-Noise 的输入包括频率、攻角、弦长、来流速度和吸力侧位移厚度，目标为经过整理的声压级（dB）。先绘制字段分布，比较均值基线、线性回归和 Ridge。声压级是对数量，不应把回归系数简单当作声压本身的线性响应。如果数据记录允许识别相关工况，讨论按工况分组对评价的影响。课程提供下载脚本，不把模拟记录冒充真实实验数据。

\scientificnote{科学项目的关键不仅是得到最小的 RMSE，还包括解释“误差对应多少个物理单位”“哪些对象或工况表现差”“模型的输入是否能够在真正预测时取得”。}

\subsection{把专业背景落实为拟合选择}
OGLE任务首先按目录说明选定区域与脉动模式，不能先把不同模式混在一起，再要求一个直线模型消除结构差异。绘制$\log_{10}(P/\mathrm{day})$与平均$I$星等后，拟合基线并检查残差。应记录目标是视星等：距离、消光、混叠及目录选择都可能影响关系；一个训练分数不能替代对这些影响的讨论。

Airfoil任务应先辨认哪些列共同定义实验工况。频率是在同一工况下可能连续扫描的变量，不能把“每一行都有不同频率”当作完全独立设备。可用攻角、弦长、流速和位移厚度的组合形成可复现的工况键，再比较按行和按工况的评价。字段精度和原始设计是否足以支持这个键，需要依据数据说明核对，不应凭少量小数任意拼组。

这两种专业实践不是把CSV路径换成另一个名字就完成。它们分别要求理解天体与观测的关系、实验条件与重复测量的关系。模型仍然使用熟悉的接口，但样本单位、变换和验证方案需要重新作出判断。

\subsection{离线机制实验与真实数据结果分开交付}
本轮全书的离线回归机制结果由明确生成的响应数据及标准离线数据得到。住房、OGLE和Airfoil提供完整获取入口、字段约定、运行脚本和分析任务；由于本次环境不能连接数据下载服务器，没有为这些外部数据填写未经运行的分数。读者获取真实数据后应运行对应脚本，生成自己的比较表，而不是把人工响应的结果当成专业观测结论。

这种区分不妨碍完成方法学习：手算与机制实验用于检验计算，真实数据项目用于检验方法在专业条件下的适用性。两者都重要，但承担不同证据角色。

\chapterpractice{标准实验完成住房数据的基线、LinearRegression、岭回归与 Lasso 比较，并用开发数据内的交叉验证选择超参数；专业实验在 OGLE 与 Airfoil 中选一项，提交字段字典、划分依据、残差图和研究范围说明。}
\chapterexercise
\begin{enumerate}
\item 试解释为何二次多项式回归仍属于线性参数模型。
\item 用同一数据比较原始单位与标准化单位的 Ridge 系数；预测是否发生必然改变？
\item 对一个有相关输入的模型，若 Lasso 将某个系数置零，能否证明该物理量不重要？
\item 比较 MAE 与 RMSE 对极端异常样本的反应，并说明该样本是否应被删除还需要什么证据。
\item 针对 OGLE 或 Airfoil，设计一个不同于随机划分的评价方案，解释它检验何种泛化能力。
\end{enumerate}

\section*{深化练习与方法检验}
本组题强调解释、计算与验证；D编号在本章内独立计数。计算题应写出所用定义和中间步骤，实验题保留实际输出，包括不支持预期的结果。
\begin{enumerate}[label=D\arabic*.]
\item 对 $x=(1,2,3)$、$y=(2,3,5)$，求含截距的最小二乘直线，计算残差及SSE，并证明残差和为零。这个性质能否用来断言不存在系统误差？
\item 考虑 $J(w)=\frac12(w-2)^2+\lambda|w|$。分别求 $\lambda=0,0.5,1,2,3$ 时的最优值，并与把 $|w|$ 改为 $w^2/2$ 的情形比较。
\item 若两列特征完全相同，线性模型能否唯一确定两个系数？构造两组不同系数给出相同预测，说明正则化与增加新信息的区别。
\item 为正的周期取 $\log_{10}P$ 后拟合直线，与对数预测值再指数还原有什么区别？在本书OGLE问题中哪个变量被取对数，目标的单位是否改变？
\end{enumerate}
